% Adapted from docs/wiki/Losses.md; edit this file, not the wiki page, for the manual.
\chapter{Losses}\label{ch:losses}

Losses are PyTorch modules accepting floating-point logits \texttt{(B,\ C,\ *spatial)}. Most accept integer
targets \texttt{(B,\ *spatial)} and \texttt{target\_format="one-hot"} for same-shaped masks or soft labels;
individual objectives document their supported target formats below.
Softmax targets should sum to one over classes; sigmoid targets should lie in \texttt{{[}0,\ 1{]}}.
Exacta does not silently normalize or clamp target probabilities.

Let \texttt{p} be activated predictions, \texttt{g} the target and \texttt{s} smoothing. Core overlap and focal losses are:

{\small\def\LTcaptype{none} % do not increment counter
\begin{longtable}[]{@{}
  >{\raggedright\arraybackslash}p{(\linewidth - 2\tabcolsep) * \real{0.2700}}
  >{\raggedright\arraybackslash}p{(\linewidth - 2\tabcolsep) * \real{0.7300}}@{}}
\toprule\noalign{}
\begin{minipage}[b]{\linewidth}\raggedright
Loss
\end{minipage} & \begin{minipage}[b]{\linewidth}\raggedright
Definition and reduction
\end{minipage} \\
\midrule\noalign{}
\endhead
\bottomrule\noalign{}
\endlastfoot
\texttt{DiceLoss} & \texttt{1\ -\ (2\ sum(p*g)+s)/(sum(p)+sum(g)+s)}, averaged across cases and selected classes \\
\texttt{DiceCELoss} & \texttt{lambda\_dice\ *\ DiceLoss\ +\ lambda\_ce\ *\ CE}; sigmoid uses BCE \\
\texttt{DiceFocalLoss} & \texttt{lambda\_dice\ *\ DiceLoss\ +\ lambda\_focal\ *\ FocalLoss} \\
\texttt{GeneralizedDiceCELoss} & \texttt{lambda\_dice\ *\ GeneralizedDiceLoss\ +\ lambda\_ce\ *\ CE}; sigmoid uses BCE \\
\texttt{TverskyLoss} & \texttt{1\ -\ (TP+s)/(TP+alpha*FP+beta*FN+s)}, averaged across cases and selected classes \\
\texttt{GeneralizedDiceLoss} & Ratio of class-weighted overlap sums; class weights derive from target volumes \\
\texttt{FocalLoss} & Focal CE/BCE with exponent \texttt{gamma}; softmax sums classes then averages voxels, sigmoid averages all entries \\
\end{longtable}
}

Overlap losses support \texttt{activation="softmax"}, \texttt{"sigmoid"} or \texttt{"none"}. Dice-CE and focal
require logits and use softmax or sigmoid. \texttt{batch\_dice=True} pools overlap statistics across
the local batch before division. The explicit global option is documented below.
\texttt{include\_background=False} excludes channel zero from overlap only. CE/BCE and focal include
all channels.

\section{Overlap equations}\label{losses:overlap-equations}

For one case and selected class, let \(p_i\) be the activated prediction and \(g_i\) the target
at included voxel \(i\). With smoothing \(s\), the soft overlap statistics are:

\[
\mathrm{TP}=\sum_i p_i g_i,\qquad
\mathrm{FP}=\sum_i p_i(1-g_i),\qquad
\mathrm{FN}=\sum_i (1-p_i)g_i.
\]

Dice and Tversky use these ratios before averaging over cases and selected classes:

\[
\mathcal{L}_{\mathrm{Dice}}
=1-\frac{2\sum_i p_i g_i+s}{\sum_i p_i+\sum_i g_i+s},\qquad
\mathcal{L}_{\mathrm{Tversky}}
=1-\frac{\mathrm{TP}+s}{\mathrm{TP}+\alpha\mathrm{FP}+\beta\mathrm{FN}+s}.
\]

Generalized Dice instead takes a single class-weighted ratio per case, where \(c\) ranges
over selected classes:

\[
\mathcal{L}_{\mathrm{GDice}}
=1-\frac{2\sum_c w_c\sum_i p_{ic}g_{ic}+s}
{\sum_c w_c\sum_i(p_{ic}+g_{ic})+s}.
\]

For present target classes, \texttt{weight\_type="square"} sets \(w_c=(\sum_i g_{ic})^{-2}\);
\texttt{"simple"} uses the inverse volume and \texttt{"uniform"} uses one. The absent-class rules below
handle zero volumes. \texttt{batch\_dice=True} includes the batch axis in the sums before taking
the ratio. Smoothing is added once per ratio, after weighting for generalized Dice.

\section{Important parameters}\label{losses:important-parameters}

The following options describe the core objectives. Squared Dice, Lovasz-hinge and
Wasserstein Dice are described at the end of this page.

Definitions differ across libraries, so select parameters intentionally:

\begin{itemize}
\tightlist
\item
  Dice/Tversky smoothing defaults to \texttt{1e-5}. Tversky's default \texttt{alpha=beta=0.5} equals Dice
  with twice the smoothing constant; alpha weights false positives and beta false negatives.
\item
  \texttt{zero\_absent=True} (Dice, Tversky and focal Tversky, since 0.1.0) gives zero loss and
  gradient to classes whose target is empty within the reduction (case, or batch with
  \texttt{batch\_dice=True}), however much is predicted; they still count in the mean. This reproduces
  segmentation\_models\_pytorch's Dice and Tversky losses with \texttt{batch\_dice=True} and \texttt{smooth=0}.
  By default, as in MONAI, an empty target still penalizes predicted probability. It runs on the
  PyTorch backend.
\item
  Generalized Dice supports \texttt{weight\_type="square"}, \texttt{"simple"} and \texttt{"uniform"}. Smoothing
  is added after class weighting. Absent classes receive the largest present-class weight.
  With inverse-volume weighting, \texttt{empty\_target="zero"} gives zero for entirely absent selected
  targets; \texttt{"uniform"} uses unit weights there. Positive smoothing is needed for a useful
  gradient on an entirely empty target. This differs from an entirely ignored case.
\item
  Focal uses \texttt{gamma=2} and \texttt{alpha=None} by default. Scalar alpha weights sigmoid positives by
  alpha and negatives by \texttt{1-alpha}; softmax assigns \texttt{1-alpha} to class zero and alpha to other
  classes. \texttt{gamma=0,\ alpha=None} reproduces CE/BCE. Softmax reduction differs from MONAI's mean
  over channels by a factor of the channel count.
\end{itemize}

\section{Backends and derivatives}\label{losses:backends-and-derivatives}

\texttt{backend="auto"} selects fused Triton for supported CUDA inputs and PyTorch otherwise.
Half/bfloat16 inputs compute in float32; float64 inputs retain their precision. Fused losses
support first derivatives. Use \texttt{backend="torch"} for higher-order derivatives; \texttt{auto} also
selects PyTorch when targets require gradients. Full-graph \texttt{torch.compile} is tested for
local losses. Distributed losses are tested in eager execution; compiled collectives are
not part of the current compatibility guarantee.

\section{Ignored labels and voxel masks}\label{losses:ignored-labels-and-voxel-masks}

Dice, Tversky, generalized Dice, focal and their regional combinations accept constructor
argument \texttt{ignore\_index} for index targets and keyword-only \texttt{mask=} in forward. Focal Tversky,
top-k CE and Lovasz-Softmax also support exclusions. Geometry/topology losses have different
requirements, described below.
Masks must be boolean, on the input device, with shape \texttt{(B,\ *spatial)} or \texttt{(B,\ 1,\ *spatial)}.
True includes a voxel; False excludes it from every class. Both exclusions combine.

\begin{lstlisting}[language=Python]
import torch
from exacta.losses import DiceCELoss

generator = torch.Generator().manual_seed(0)
logits = torch.randn(2, 3, 4, 4, 4, generator=generator, requires_grad=True)
target = torch.randint(3, (2, 4, 4, 4), generator=generator)
target[:, 0] = -100
mask = torch.ones_like(target, dtype=torch.bool)
mask[:, :, 0] = False
loss = DiceCELoss(ignore_index=-100)(logits, target, mask=mask)
loss.backward()
assert torch.count_nonzero(logits.grad[:, :, 0]) == 0
\end{lstlisting}

CE/focal normalize over included voxels (and channels for sigmoid). Overlap averages exclude
entirely ignored cases; a fully ignored batch returns zero loss and gradients even with zero
smoothing. Annotated background voxels remain valid observations. Ignored invalid labels and
nonfinite values are excluded before activation/reduction. Valid out-of-range labels still
raise eagerly; compiled Triton returns NaN. Soft/one-hot targets use \texttt{mask}, not \texttt{ignore\_index}.

The GPU kernels read exclusion information directly and count valid voxels in their existing
passes. They do not create full-volume masked probability or one-hot buffers.

\section{Loss combinations and CE weights}\label{losses:loss-combinations-and-ce-weights}

\texttt{DiceFocalLoss} adds Dice and focal terms; \texttt{GeneralizedDiceCELoss} adds generalized Dice and
CE/BCE. Their GPU implementations share one activation/target pass forward and one backward.
The corresponding native PyTorch compositions provide a reference and higher-order gradients.
Coefficients must be finite and nonnegative, with at least one positive term. A coefficient
of zero skips that term, including otherwise undefined zero-smoothing overlap ratios.

\begin{lstlisting}[language=Python]
from exacta.losses import DiceFocalLoss, GeneralizedDiceCELoss

dice_focal = DiceFocalLoss(include_background=False, gamma=2, alpha=0.25,
                          lambda_dice=1, lambda_focal=0.5)
gdl_ce = GeneralizedDiceCELoss(include_background=False, empty_target="uniform",
                              class_weight=[0.2, 1.0, 2.0])
\end{lstlisting}

\texttt{DiceCELoss} and \texttt{GeneralizedDiceCELoss} accept \texttt{class\_weight}, one nonnegative finite value
per input channel, stored as a module buffer. These weights apply only to CE/BCE. Normalization
matches \href{https://docs.pytorch.org/docs/2.14/generated/torch.nn.CrossEntropyLoss.html}{PyTorch's cross-entropy definition}
and depends on target format:

{\small\def\LTcaptype{none} % do not increment counter
\begin{longtable}[]{@{}
  >{\raggedright\arraybackslash}p{(\linewidth - 4\tabcolsep) * \real{0.3000}}
  >{\raggedright\arraybackslash}p{(\linewidth - 4\tabcolsep) * \real{0.4500}}
  >{\raggedright\arraybackslash}p{(\linewidth - 4\tabcolsep) * \real{0.2500}}@{}}
\toprule\noalign{}
\begin{minipage}[b]{\linewidth}\raggedright
Pointwise term
\end{minipage} & \begin{minipage}[b]{\linewidth}\raggedright
Weighted numerator
\end{minipage} & \begin{minipage}[b]{\linewidth}\raggedright
Denominator
\end{minipage} \\
\midrule\noalign{}
\endhead
\bottomrule\noalign{}
\endlastfoot
Softmax, index target & Sum of \texttt{-w{[}y{]}\ log\ p{[}y{]}} over included voxels & Sum of included \texttt{w{[}y{]}} \\
Softmax, soft/one-hot target & Sum of \texttt{-w{[}c{]}\ g{[}c{]}\ log\ p{[}c{]}} & Number of included voxels \\
Sigmoid BCE & Class weight multiplies both positive and negative terms & Included voxels × channels \\
\end{longtable}
}

Zero total target weight returns zero CE and zero CE gradient. Exact one-hot and index labels
therefore produce different weighted CE scales, as in PyTorch. \texttt{include\_background=False}
still affects only the overlap term. Use focal's \texttt{alpha} for its documented positive/negative
or background/foreground weighting.

\section{Deep supervision}\label{losses:deep-supervision}

\texttt{DeepSupervisionLoss} applies one scalar loss to explicitly paired resolutions. Weights are
used as supplied, without normalization; zero-weight levels are not evaluated and receive
no gradient. Targets and optional masks must already match each prediction's spatial shape.
The wrapper does not resize labels or choose downsampling conventions.

\begin{lstlisting}[language=Python]
from exacta.losses import DeepSupervisionLoss, DiceCELoss

loss_fn = DeepSupervisionLoss(DiceCELoss(ignore_index=-100), weights=[1.0, 0.5, 0.25])
# Sequences of three tensors, ordered consistently from fine to coarse:
loss = loss_fn(predictions, targets, masks=voxel_masks)
\end{lstlisting}

\section{Global distributed losses}\label{losses:global-distributed-losses}

Set \texttt{distributed=True} to pool statistics over the default \texttt{torch.distributed} process group.
Overlap losses and their combinations also require \texttt{batch\_dice=True}. This computes the
objective for the concatenated global batch, including generalized-Dice class weights and
weighted CE denominators. It differs from averaging independently computed rank-local Dice.

\begin{lstlisting}[language=Python]
# After initializing the process group and wrapping the model in ordinary DDP:
loss_fn = GeneralizedDiceCELoss(batch_dice=True, distributed=True,
                               include_background=False, ignore_index=-100)
loss = loss_fn(ddp_model(images), labels, mask=annotated)
loss.backward()
\end{lstlisting}

Every rank must use identical loss configuration and backend, and call forward and backward
in the same order, including ranks with zero cases or entirely ignored data. An entirely
ignored or empty global batch returns zero. Do not multiply the loss by local batch size or
divide by world size: the differentiable statistics sum and DDP's gradient averaging together
produce concatenated-batch \textbf{parameter gradients}; see
\href{https://docs.pytorch.org/docs/2.14/generated/torch.nn.parallel.DistributedDataParallel.html}{DDP's gradient-averaging convention}.
Logit/target gradients before DDP averaging
carry a world-size factor; custom gradient communication hooks need their own scaling analysis.
Gradient accumulation across microbatches averages microbatch objectives; it does not pool
their nonlinear Dice ratios into one larger-batch ratio.

Only small statistics communicate. NCCL keeps these on the rank's CUDA device. Gloo uses CPU
buffers, also allowing two CUDA workers on a single GPU for correctness tests. Local tests
cover CPU and fused CUDA with two Gloo ranks, uneven batches, empty ranks, ignored ranks and
all-ignored batches. \textbf{The multi-GPU NCCL path is experimental and unverified in 0.1.2:}
the two-GPU tests have not been run on suitable hardware. Gloo checks do not establish NCCL
correctness or multi-GPU training throughput. Loss configuration and valid input shapes
must agree across ranks; as with ordinary distributed training, a rank-local error can interrupt its peers.

\section{Boundary, topology and hard-example objectives}\label{losses:boundary-topology-and-hard-example-objectives}

Additional objectives include \texttt{FocalTverskyLoss}, \texttt{TopKCELoss}, \texttt{DiceTopKCELoss}, \texttt{LovaszSoftmaxLoss},
\texttt{BoundaryLoss}, \texttt{HausdorffDTLoss}, \texttt{SoftClDiceLoss} and \texttt{SkeletonRecallLoss}. Their formulas, reductions, defaults,
reference conventions and CPU/GPU paths are documented in the
\href{https://github.com/rodolfo-viana/exacta/wiki/Advanced-segmentation}{advanced segmentation guide}.

Focal Tversky supports the global distributed reduction above. Top-k/Lovasz selection and
geometry/topology losses operate on each rank's local batch. Geometry masks affect reduction
only; complete target geometry is required. Geometry/topology losses have no \texttt{ignore\_index},
and soft-clDice has no mask argument. Distance-transform and skeleton objectives should be
combined with regional supervision according to the task.

\section{Soft labels}\label{losses:soft-labels}

With soft targets \(g_i\in(0,1)\), the product overlap \(\sum_i p_i g_i\) makes Dice and Tversky
losses minimal away from \(p=g\): they push predictions towards 0 and 1. Label smoothing,
distillation and averaged annotations therefore need another form. \texttt{soft\_label=True} in
\texttt{DiceLoss}, \texttt{TverskyLoss}, \texttt{FocalTverskyLoss}, \texttt{GeneralizedDiceLoss}, \texttt{DiceCELoss},
\texttt{DiceFocalLoss} and \texttt{GeneralizedDiceCELoss} uses the semimetric statistics of
\href{https://doi.org/10.1007/978-3-031-43898-1_46}{Wang et al.~(MICCAI 2023)}:

\[
\mathrm{TP}=\tfrac{1}{2}\Big(\sum_i p_i+\sum_i g_i-\sum_i\lvert p_i-g_i\rvert\Big),\qquad
\mathrm{FP}=\sum_i p_i-\mathrm{TP},\qquad \mathrm{FN}=\sum_i g_i-\mathrm{TP}.
\]

For binary targets \(\lvert p_i-g_i\rvert=p_i+g_i-2p_ig_i\), so values equal the ordinary losses
exactly, and so do gradients except at voxels where \(p_i=g_i\in\{0,1\}\), where the zero
subgradient of \(\lvert p_i-g_i\rvert\) applies; through softmax or sigmoid, whose derivative
vanishes there, logit gradients agree. For soft targets every ratio is maximal only at \(p=g\) (for Tversky,
when both weights are positive). It matches MONAI's \texttt{soft\_label=True}. The option acts on
one-hot targets and runs on the PyTorch backend; index targets are hard and keep the fused
path. CE terms are
unchanged: cross-entropy is already minimal at \(p=g\). \texttt{squared\_pred=True} is too
(\(\sum p^2+\sum g^2-2\sum pg=\sum(p-g)^2\)), so the two options are exclusive.

\begin{lstlisting}[language=Python]
from exacta.losses import DiceCELoss

loss_fn = DiceCELoss(target_format="one-hot", soft_label=True)
loss = loss_fn(logits, smoothed_one_hot_targets)
\end{lstlisting}

\section{Squared Dice denominator}\label{losses:squared-dice-denominator}

\texttt{squared\_pred=True} is available in \texttt{DiceLoss}, \texttt{DiceCELoss}, \texttt{DiceFocalLoss}
and \texttt{DiceTopKCELoss}. Both prediction and target terms in the denominator are squared:

\[
\mathcal{L}_{\mathrm{squared\ Dice}}
=1-\frac{2\sum_i p_i g_i+s}{\sum_i p_i^2+\sum_i g_i^2+s}.
\]

The numerator retains the product overlap, including for soft targets. With binary targets,
\(g_i^2=g_i\). The ordinary denominator remains the default. Smoothing, background exclusion,
masking and case/batch reduction follow ordinary Dice. The option is named consistently with
\href{https://monai.readthedocs.io/en/stable/losses.html\#diceloss}{MONAI's squared-prediction option};
the equation above defines Exacta's soft-target convention explicitly.

\begin{lstlisting}[language=Python]
from exacta.losses import DiceCELoss

loss_fn = DiceCELoss(squared_pred=True, include_background=False)
\end{lstlisting}

PyTorch supports differentiation through predictions and soft targets. The fused CUDA path
includes the squared denominator in both forward statistics and analytic backward, and
supports global distributed Dice/CE reductions. Fused higher-order derivatives remain
unsupported; use \texttt{backend="torch"} for those. Dice-top-k still selects hard examples locally.

\section{Binary and multilabel Lovasz-hinge}\label{losses:binary-and-multilabel-lovasz-hinge}

\texttt{LovaszHingeLoss} follows the binary real-valued-score formulation in the
\href{https://github.com/bermanmaxim/LovaszSoftmax/blob/master/pytorch/lovasz_losses.py}{authors' implementation}.
Given binary labels \(y_i\) and raw logits \(z_i\), margin errors are
\(e_i=1-z_i(2y_i-1)\). Sort the errors decreasingly, and let \(\Delta J_j\) be successive
increments of the Jaccard error for that ordering. The objective is:

\[
\mathcal{L}_{\mathrm{Lovasz\ hinge}}=\sum_j\max(e_{(j)},0)\,\Delta J_j.
\]

Binary input has shape \texttt{(B,1,*spatial)} with target \texttt{(B,*spatial)}. Independent overlapping
labels use \texttt{multilabel=True} and a same-shaped binary target. No sigmoid is applied and
soft targets are rejected. Masks can be shared or per-label; \texttt{ignore\_index} applies only
to integer binary targets. Excluded invalid logits/targets do not affect the result.

\texttt{per\_image=True} (default) averages observed case/label losses. False pools the local batch
per label before averaging labels. An absent positive class still penalizes false positives;
unobserved labels are omitted, and entirely excluded inputs give differentiable zero.
Stable native PyTorch sorting uses flattened voxel order for ties. The subgradient at a
zero hinge margin is zero. This is a local objective under DDP, without global sorting.

\begin{lstlisting}[language=Python]
from exacta.losses import LovaszHingeLoss

binary_loss = LovaszHingeLoss(ignore_index=-1)
region_loss = LovaszHingeLoss(multilabel=True, per_image=True)
\end{lstlisting}

These losses support both 2D and 3D medical segmentation. They are alternative
objectives with different gradients, not a claim of improved accuracy for every task.

\section{Generalized Wasserstein Dice}\label{losses:generalized-wasserstein-dice}

\texttt{GeneralizedWassersteinDiceLoss} adds explicit costs for confusing mutually exclusive labels.
It follows the crisp-target formulation of \href{https://arxiv.org/abs/1707.00478}{Fidon et al.~(2017)},
with the alternative \texttt{GDL} weighting described by
\href{https://monai.readthedocs.io/en/stable/losses.html\#generalizedwassersteindiceloss}{MONAI}.
It supports 2D images and 3D volumes. Costs describe relationships between labels; voxel
spacing and physical boundary distances do not enter this objective.

Let \(M\) be the supplied class-cost matrix divided by its positive maximum, \(y_i\) the integer
target at observed voxel \(i\), and \(p_{ic}\) the softmax probability for class \(c\). The expected
cost against that crisp target is

\[
w_i=\sum_c M_{y_i,c}\,p_{ic}.
\]

For each case, define \(N_c=\sum_i\mathbf{1}[y_i=c]\) over the observed domain and
\(T=\sum_i\alpha_{y_i}(1-w_i)\). The two modes use different weights and error sums:

{\small\def\LTcaptype{none} % do not increment counter
\begin{longtable}[]{@{}lll@{}}
\toprule\noalign{}
\texttt{weighting\_mode} & Class weights \(\alpha_c\) & Error \(E\) \\
\midrule\noalign{}
\endhead
\bottomrule\noalign{}
\endlastfoot
\texttt{"default"} & \(\alpha_0=0\); \(\alpha_c=1\) for \(c>0\) & \(\sum_i w_i\) \\
\texttt{"GDL"} & \(1/(N_c+1)\), including label zero & \(\sum_i\alpha_{y_i}w_i\) \\
\end{longtable}
}

Both return the following loss, with the same smoothing constant \(s\) in both terms:

\[
\mathcal{L}=1-\frac{2T+s}{2T+E+s}.
\]

The constructor requires a finite, nonnegative, symmetric square matrix with zero diagonal
and a positive maximum. \texttt{default} additionally requires maximal costs between background
(label zero) and every foreground class, matching the assumption behind the simplified
foreground formula. \texttt{GDL} permits other background costs. Triangle inequalities are not
checked; a zero off-diagonal cost makes that label confusion cost-free. Specify these costs
from the task's label relationships, and assess any tuning using training/validation data.

The following synthetic example gives confusion between foreground labels a smaller cost
than confusion with background. The same loss accepts \texttt{(B,C,H,W)} or \texttt{(B,C,D,H,W)} logits.

\begin{lstlisting}[language=Python]
import torch
from exacta.losses import GeneralizedWassersteinDiceLoss

device = torch.device("cuda" if torch.cuda.is_available() else "cpu")
rng = torch.Generator().manual_seed(879)
logits = torch.randn((2, 3, 16, 20), generator=rng).to(device).requires_grad_()
target = torch.randint(3, (2, 16, 20), generator=rng).to(device)
target[0, :2] = -1

loss_fn = GeneralizedWassersteinDiceLoss(
    [[0, 1, 1], [1, 0, 0.25], [1, 0.25, 0]],
    weighting_mode="default", ignore_index=-1,
).to(device=device, dtype=logits.dtype)
loss = loss_fn(logits, target)
loss.backward()
assert torch.isfinite(loss) and torch.isfinite(logits.grad).all()
\end{lstlisting}

The default \texttt{reduction="mean"} averages observed cases equally. \texttt{"sum"} sums case losses;
\texttt{"none"} returns \texttt{(B,)}. Fully excluded cases and empty inputs return differentiable zero,
and are omitted from the mean. A zero denominator with \texttt{smooth=0} also returns zero.
With default weights and \texttt{smooth=0}, an observed background-only case with any nonzero error
has loss one and no useful foreground Dice gradient; positive smoothing or the \texttt{GDL} mode
changes that behavior. This differs from an entirely unobserved case.

This API takes integer targets and logits; it has no sigmoid/multilabel, soft-target or
cross-rank pooling option. It supports shared boolean masks, ignored indices, half/bfloat16
computation in float32, float64, higher derivatives and full-graph \texttt{torch.compile} on CPU/CUDA.
Eager calls reject out-of-range observed labels. Compiled calls return NaN for invalid labels
without unsafe indexing. Excluded labels/nonfinite logits are sanitized before activation.

The implementation uses native PyTorch gather/reduction kernels and an integer class histogram,
without expanding targets to one-hot. Move the loss to the logits' device/dtype once and
benchmark \texttt{torch.compile(loss\_fn)} with forward \textbf{and backward} on your workload. The cost
matrix is a detached normalized buffer, included in \texttt{state\_dict}; it is not a learnable parameter.
